\documentclass[11pt]{amsart}

\usepackage[utf8]{inputenc}
\usepackage[T1]{fontenc}
\usepackage{lmodern}
\usepackage{amsmath,amssymb,amsthm,mathtools}
\usepackage[margin=1in]{geometry}
\usepackage{microtype}
\usepackage{enumitem}
\usepackage{xcolor}
\usepackage[colorlinks=true,
  linkcolor=blue!45!black,
  citecolor=blue!45!black,
  urlcolor=blue!45!black]{hyperref}

\newtheorem{theorem}{Theorem}[section]
\newtheorem{proposition}[theorem]{Proposition}
\newtheorem{lemma}[theorem]{Lemma}
\newtheorem{corollary}[theorem]{Corollary}
\theoremstyle{definition}
\newtheorem{definition}[theorem]{Definition}
\newtheorem{problem}[theorem]{Problem}
\newtheorem{hypothesis}[theorem]{Hypothesis}
\theoremstyle{remark}
\newtheorem{remark}[theorem]{Remark}
\newtheorem{warning}[theorem]{Warning}

\DeclareMathOperator{\Dom}{Dom}
\DeclareMathOperator{\Ran}{Ran}
\DeclareMathOperator{\spanop}{span}
\DeclareMathOperator{\dist}{dist}
\DeclareMathOperator{\Repart}{Re}
\DeclareMathOperator{\Impart}{Im}
\newcommand{\C}{\mathbb C}
\newcommand{\R}{\mathbb R}
\newcommand{\Hh}{\mathbb H}
\newcommand{\Bcal}{\mathcal B}
\newcommand{\Kcal}{\mathcal K}
\newcommand{\Hcal}{\mathcal H}
\newcommand{\Dcal}{\mathcal D}
\newcommand{\ip}[2]{\left\langle #1,#2\right\rangle}
\newcommand{\norm}[1]{\left\lVert #1\right\rVert}
\newcommand{\abs}[1]{\left\lvert #1\right\rvert}
\newcommand{\rhoB}{\rho}
\newcommand{\eps}{\varepsilon}

\title{\bfseries The Lax--Phillips--Pavlov--Faddeev\\
Asymptotic Criterion for the Riemann Hypothesis\\[4pt]
\large A convention-fixed problem specification and solver handoff}
\author{OpenAI Codex}
\date{10 August 2026}

\begin{document}
\maketitle

\begin{abstract}
This document fixes a mathematically unambiguous target corresponding to the
Pavlov--Faddeev asymptotic criterion quoted in later accounts of the
Lax--Phillips scattering theory of the automorphic wave equation.  It defines
the reduced positive energy space, the interaction semigroup, its generator
and resolvent, and the precise regularized operator norm.  In the convention
adopted here, the quoted RH-equivalence is
\[
 \mathrm{RH}
 \quad\Longleftrightarrow\quad
 \limsup_{t\to\infty}\frac1t
 \log\norm{Z(t)(B+iI)^{-1}}_{E\to E}=-\frac14.
\]
The resolvent factor is essential.  Equivalently, the estimate measures the
operator norm of $Z(t)$ from the generator domain with its graph norm into the
ordinary energy space.  The unregularized norm $\norm{Z(t)}_{E\to E}$ and the
fixed-vector quantity $\norm{Z(t)f}_{E}$ are different objects and are recorded
separately.  Conditional on the physical-to-model eigenvalue correspondence,
the implication from the regularized upper bound to RH is proved below.  The
reverse implication is attributed to Pavlov--Faddeev by a 2017 survey but is
not reconstructed from the 1972/1975 primary paper here.  Accordingly, this
handoff separates three tasks: verifying the historical theorem, reproving its
conditional analytic direction, and deciding the regularized estimate
unconditionally.  Only the last task would decide RH.
\end{abstract}

\tableofcontents

\section{The target statement in one box}

The notation in this box is defined in full in Sections~\ref{sec:energy}--\ref{sec:regularized}.

\begin{center}
\fbox{\begin{minipage}{0.92\textwidth}
Let $(\Kcal,\ip{\cdot}{\cdot}_{E})$ be the reduced Lax--Phillips
interaction Hilbert space for the modular automorphic wave equation.  Let
$Z=(Z(t))_{t\geq0}$ be its interaction semigroup and let $B$ be the ordinary
$C_{0}$-semigroup generator, so that $Z(t)=e^{tB}$.  Assume the explicitly
listed energy-space, scattering, resolvent, and zeta-eigenvalue identification
hypotheses in this document.  In particular, assume $-i\in\rho(B)$ and put
\[
 M(t):=\norm{Z(t)(B+iI)^{-1}}_{E\to E},
 \qquad
 \omega_{\mathrm{reg}}(B)
 :=\limsup_{t\to\infty}\frac1t\log M(t).
\]
The mathematical target is $\omega_{\mathrm{reg}}(B)=-1/4$.  The historical
equivalence quoted by Takhtajan et al.\ is
\[
 \boxed{\quad
 \mathrm{RH}\quad\Longleftrightarrow\quad
 \omega_{\mathrm{reg}}(B)=-\frac14.
 \quad}
\]
Under the zeta-eigenvalue identification,
$\omega_{\mathrm{reg}}(B)\geq-1/4$ unconditionally, so the right-hand
condition can equivalently be written
\[
 \mathrm{RH}\quad\Longleftrightarrow\quad
 \omega_{\mathrm{reg}}(B)\leq-\frac14.
\]
The norm is
\[
 M(t)=
 \sup_{0\neq h\in\Kcal}
 \frac{\norm{Z(t)(B+iI)^{-1}h}_{E}}{\norm{h}_{E}}
 =
 \sup_{0\neq f\in\Dom B}
 \frac{\norm{Z(t)f}_{E}}{\norm{(B+iI)f}_{E}}.
\]
\end{minipage}}
\end{center}

The last equality is the most important norm specification in this document:
the source norm is a graph norm on $\Dom B$, while the target norm is the
ordinary positive energy norm.

\section{Arithmetic normalization}

\begin{definition}[Zeta and completed zeta]
For $\Repart s>1$ let
\[
 \zeta(s)=\sum_{n=1}^{\infty}n^{-s}.
\]
Its meromorphic continuation has a simple pole at $s=1$.  Define
\[
 \xi(s):=\frac12s(s-1)\pi^{-s/2}\Gamma\!\left(\frac{s}{2}\right)\zeta(s).
\]
Then $\xi$ is entire and satisfies
\[
 \xi(s)=\xi(1-s)=\overline{\xi(\bar s)}.
\]
A \emph{nontrivial zero} is a zero $\rho=\beta+i\gamma$ of $\xi$.
\end{definition}

\begin{definition}[Riemann hypothesis]
The Riemann hypothesis is the assertion
\[
 \mathrm{RH}:\qquad \xi(\rho)=0\ \Longrightarrow\ \Repart\rho=\frac12.
\]
The multiset of zeros is invariant under $\rho\mapsto\bar\rho$ and
$\rho\mapsto1-\rho$.
\end{definition}

The factor $1/4$ in the target asymptotic is forced by the following
model normalization, once the physical-to-model identification in
Hypothesis~\ref{hyp:model} is available:
\[
 a=\frac{\rho}{2}
 \quad\longmapsto\quad
 \lambda_{\rho}=-\bar a=-\frac{\bar\rho}{2},
 \qquad
 \Repart\lambda_{\rho}=-\frac{\beta}{2}.
\]
Thus the critical line $\beta=1/2$ maps to the semigroup spectral line
$\Repart\lambda=-1/4$.

\section{The automorphic wave energy space}\label{sec:energy}

\subsection{The spatial operator and wave equation}

Let
\[
 X=\operatorname{PSL}_{2}(\mathbb Z)\backslash\Hh,
 \qquad
 d\mu(z)=\frac{dx\,dy}{y^{2}},\qquad z=x+iy,
\]
and let
\[
 \Delta_X=-y^{2}(\partial_x^{2}+\partial_y^{2})
\]
denote the nonnegative self-adjoint Friedrichs realization of the automorphic
Laplacian in $L^{2}(X,d\mu)$.  Following the Lax--Phillips normalization, set
\[
 A:=\Delta_X-\frac14I.
\]
The corresponding automorphic wave equation is
\begin{equation}\label{eq:wave}
 u_{tt}+Au=0.
\end{equation}
The absolutely continuous spectrum of $A$ starts at $0$.  Below-threshold
modes and threshold null directions are responsible for indefiniteness or
degeneracy of the raw energy form.  Positive-energy cusp forms are removed for
a different reason: they are point-spectrum standing waves rather than
scattering states.

\subsection{Energy form and positive reduction}

Let $P_{\mathrm{ac}}$ be the absolutely continuous spectral projection of
$A$, put
\[
 \mathfrak h_{c}:=\Ran P_{\mathrm{ac}},
 \qquad
 A_{c}:=A\big|_{\mathfrak h_{c}},
\]
and note that $A_{c}\geq0$.  For sufficiently regular Cauchy data
$f=(f_{0},f_{1})$ and $g=(g_{0},g_{1})$ in this sector, define
\begin{equation}\label{eq:energy-form}
 E(f,g)
 :=\ip{A_c^{1/2}f_{0}}{A_c^{1/2}g_{0}}_{L^{2}(X)}
   +\ip{f_{1}}{g_{1}}_{L^{2}(X)},
\end{equation}
Equivalently, on a common operator core,
\[
 E(f,g)=\ip{A_cf_{0}}{g_{0}}_{L^{2}(X)}
       +\ip{f_{1}}{g_{1}}_{L^{2}(X)}.
\]

Define
\[
 \dot{\mathfrak h}^{\,1}_{A_c}
 :=
 \overline{\Dom(A_c^{1/2})/\ker A_c}^{\,
            \norm{A_c^{1/2}\,\cdot}_{2}}
\]
and set
\begin{equation}\label{eq:modern-energy-space}
 \Hcal_{E}^{c}
 :=\dot{\mathfrak h}^{\,1}_{A_c}\oplus\mathfrak h_c,
 \qquad
 \norm{(f_{0},f_{1})}_{E}^{2}
 :=\norm{A_{c}^{1/2}f_{0}}_{2}^{2}+\norm{f_{1}}_{2}^{2}.
\end{equation}
Thus
\[
 \norm{f}_{E}:=E(f,f)^{1/2},
 \qquad f\in\Hcal_{E}^{c}.
\]
This spectral construction explicitly removes the point spectrum.  It is a
modern convention-fixed realization of the no-standing-waves energy space.
The historical construction starts from an auxiliary positive completion of
an indefinite energy form and then passes through primed and double-primed
quotients/subspaces.  The identification of those historical objects with
\eqref{eq:modern-energy-space} is an input, not a consequence of notation.

\begin{hypothesis}[Energy-space identification]\label{hyp:energy}
There is an energy-unitary identification between
\eqref{eq:modern-energy-space} and the final positive continuous
Pavlov--Faddeev/Lax--Phillips energy space on which their interaction semigroup
and regularized estimate are defined.
\end{hypothesis}

\begin{remark}[Norm convention]
Throughout this document $\norm{f}_{E}$ is the square root of the conserved
quadratic energy, as in Lax's Hilbert-space convention.  Every operator norm is
induced by this Hilbert norm.
\end{remark}

\subsection{The wave group}

Functional calculus gives a unitary group $U(t)$ on $\Hcal_{E}^{c}$:
\[
 U(t)\binom{f_{0}}{f_{1}}
 =\binom{
   \cos(tA_{c}^{1/2})f_{0}
   +A_{c}^{-1/2}\sin(tA_{c}^{1/2})f_{1}}
 {-A_{c}^{1/2}\sin(tA_{c}^{1/2})f_{0}
   +\cos(tA_{c}^{1/2})f_{1}}.
\]
The formula is interpreted by spectral calculus at the threshold.  It satisfies
\[
 U(t+s)=U(t)U(s),\qquad U(t)^{*}=U(-t),\qquad
 \norm{U(t)f}_{E}=\norm{f}_{E}.
\]

\section{Incoming, outgoing, and interaction spaces}\label{sec:interaction}

Let $\Dcal_{-},\Dcal_{+}\subset\Hcal_{E}^{c}$ denote the final
double-primed (causal, modified) incoming and outgoing subspaces of the
automorphic Lax--Phillips construction.  Rather than hide their special
properties in an unstated convention, we isolate exactly what is used.

\begin{hypothesis}[Orthogonal automorphic scattering system]\label{hyp:LP}
The subspaces $\Dcal_{\pm}$ are closed and satisfy
\begin{align*}
 U(t)\Dcal_{+}&\subset\Dcal_{+} &&(t\geq0),
 &U(t)\Dcal_{-}&\subset\Dcal_{-} &&(t\leq0),\\
 \bigcap_{t\in\R}U(t)\Dcal_{\pm}&=\{0\},
 &\overline{\bigcup_{t\in\R}U(t)\Dcal_{\pm}}&=\Hcal_{E}^{c},\\
 \Dcal_{-}&\perp_{E}\Dcal_{+}.&&
\end{align*}
Moreover, the operators defined in \eqref{eq:LP-semigroup} below map
$\Kcal$ into itself, form a $C_{0}$ contraction semigroup, and obey the
compression identity following that display.
\end{hypothesis}

The adjective modified records the causal correction and removal of the
exceptional/standing-wave directions.  These orthogonality and compression
properties are special to the final automorphic construction; they are not
asserted for arbitrary incoming/outgoing subspaces.

Let $P_{\Dcal_{\pm}}^{E}$ be the energy-orthogonal projections and set
\[
 Q_{\pm}:=I-P_{\Dcal_{\pm}}^{E},
 \qquad
 \Kcal:=\Hcal_{E}^{c}\ominus_{E}(\Dcal_{-}\oplus\Dcal_{+}).
\]
The Lax--Phillips interaction semigroup is
\begin{equation}\label{eq:LP-semigroup}
 Z(t):=Q_{+}U(t)Q_{-}\big|_{\Kcal},\qquad t\geq0.
\end{equation}
Under Hypothesis~\ref{hyp:LP},
\[
 Z(t)=P_{\Kcal}^{E}U(t)\big|_{\Kcal}.
\]
The family $Z=(Z(t))_{t\geq0}$ is then a strongly continuous contraction
semigroup on $(\Kcal,\norm{\cdot}_{E})$.

Some sources denote the final reduced objects by $\Hcal'_{c}$, $\Kcal''$,
$\Dcal_{\pm}''$, and $Z''(t)$.  In this document those double primes are
suppressed: every occurrence of $\Kcal$ or $Z(t)$ refers to the final reduced
positive-energy interaction space and semigroup, not to the raw indefinite
space.

\section{Generator and resolvent}\label{sec:generator}

\begin{definition}[Ordinary semigroup generator]
The generator of $Z$ is the closed densely defined operator
\[
 Bf:=\lim_{h\downarrow0}\frac{Z(h)f-f}{h},
 \]
with domain
\[
 \Dom B:=\left\{f\in\Kcal:
 \lim_{h\downarrow0}\frac{Z(h)f-f}{h}
 \text{ exists in }\norm{\cdot}_{E}\right\}.
\]
Thus $Z(t)=e^{tB}$ in the $C_{0}$-semigroup sense.
\end{definition}

For $\lambda\in\rho(B)$, write
\[
 R(\lambda,B):=(\lambda I-B)^{-1}\in\Bcal(\Kcal).
\]
Because $Z$ is contractive, the Hille--Yosida theorem guarantees
$\{\Repart\lambda>0\}\subset\rho(B)$; in particular $R(1,B)$ always exists.
The historical criterion instead uses
$(B+iI)^{-1}=-R(-i,B)$, whose minus sign is irrelevant to the norm.

\begin{hypothesis}[Historical resolvent point]\label{hyp:minus-i}
For the convention-fixed physical generator $B$, one has $-i\in\rho(B)$.
\end{hypothesis}

Hypothesis~\ref{hyp:minus-i} follows from the full model/resolvent theorem if
that theorem gives the spectrum stated in Hypothesis~\ref{hyp:model}.  Until
that identification is verified, $R(1,B)$ provides a guaranteed
convention-free regularizer; Lemma~\ref{lem:respoint} shows that the exponent
is the same whenever both resolvents exist.

\begin{remark}[Alternative generator convention]
If another source writes $Z(t)=e^{-itG}$, then the generator used here is
$B=-iG$.  One must transform the resolvent factor at the same time.  Statements
about $B+iI$ cannot be copied between these conventions without this
conversion.
\end{remark}

\section{The regularized norm}\label{sec:regularized}

\begin{definition}[Energy operator norm]
For $T\in\Bcal(\Kcal)$,
\[
 \norm{T}_{E\to E}
 :=\sup_{0\neq f\in\Kcal}
   \frac{\norm{Tf}_{E}}{\norm{f}_{E}}.
\]
This is an operator norm on the reduced interaction Hilbert space.  It is not
the value of the energy quadratic form and not a fixed-vector norm.
\end{definition}

\begin{definition}[Regularized Lax--Phillips quantity]
For $t\geq0$ define
\begin{align}
 M(t)&:=\norm{Z(t)(B+iI)^{-1}}_{E\to E},
 \label{eq:M}\\
 \omega_{\mathrm{reg}}(B)
 &:=\limsup_{t\to\infty}\frac1t\log M(t).
 \label{eq:omega-reg}
\end{align}
The order of the two factors is immaterial because a $C_{0}$-semigroup
commutes with every resolvent of its generator:
\[
 Z(t)(B+iI)^{-1}=(B+iI)^{-1}Z(t).
\]
\end{definition}

\begin{proposition}[Exact graph-norm interpretation]\label{prop:graph}
Define
\[
 \norm{f}_{B,i}:=\norm{(B+iI)f}_{E},\qquad f\in\Dom B.
\]
Then $\norm{\cdot}_{B,i}$ is equivalent to the usual graph norm
$\norm{f}_{E}+\norm{Bf}_{E}$, and
\begin{equation}\label{eq:graph-equality}
 M(t)=
 \norm{Z(t)}_{(\Dom B,\norm{\cdot}_{B,i})\to(\Kcal,\norm{\cdot}_{E})}
 =\sup_{0\neq f\in\Dom B}
   \frac{\norm{Z(t)f}_{E}}{\norm{(B+iI)f}_{E}}.
\end{equation}
\end{proposition}

\begin{proof}
The map $B+iI:(\Dom B,\norm{\cdot}_{B,i})\to\Kcal$ is an isometric
bijection.  Substituting $h=(B+iI)f$ in \eqref{eq:M} proves
\eqref{eq:graph-equality}.  The upper comparison with the standard graph norm
is immediate.  Conversely, for $h=(B+iI)f$,
\[
 f=(B+iI)^{-1}h,
 \qquad
 Bf=h-if,
\]
so $\norm{f}_{E}+\norm{Bf}_{E}\leq
(1+2\norm{(B+iI)^{-1}})\norm{h}_{E}$.  Hence the norms are equivalent.
\end{proof}

\begin{lemma}[The resolvent point is inessential]\label{lem:respoint}
For fixed $\lambda,\mu\in\rho(B)$, the functions
\[
 M_{\lambda}(t):=\norm{Z(t)R(\lambda,B)}_{E\to E},
 \qquad
 M_{\mu}(t):=\norm{Z(t)R(\mu,B)}_{E\to E}
\]
differ by multiplicative constants independent of $t$.  Consequently their
logarithmic limsup exponents are equal.
\end{lemma}

\begin{proof}
The operator
\[
 C_{\lambda\mu}:=(\lambda I-B)(\mu I-B)^{-1}
 =I+(\lambda-\mu)R(\mu,B)
\]
is boundedly invertible.  The identity
$R(\mu,B)=R(\lambda,B)C_{\lambda\mu}$ and commutation with $Z(t)$ give
\[
 M_{\mu}(t)\leq\norm{C_{\lambda\mu}}M_{\lambda}(t).
\]
Interchanging $\lambda$ and $\mu$ gives the reverse comparison.
\end{proof}

\begin{remark}
The regularizer is therefore not the special number $i$; it is the insertion
of one fixed generator resolvent, equivalently one graph norm of regularity.
It leaves the output energy norm unchanged.
\end{remark}

\section{The quoted RH-equivalent criterion and its logical status}
\label{sec:criterion}

\begin{definition}[Regularized asymptotic statement]\label{def:PF}
For the convention-fixed modular automorphic Lax--Phillips semigroup, let
$\mathsf{PF}_{\mathrm{reg}}$ denote the assertion
\begin{equation}\label{eq:PF-criterion}
 \boxed{\quad
 \limsup_{t\to\infty}\frac1t
 \log\norm{Z(t)(B+iI)^{-1}}_{E\to E}
 =-\frac14.
 \quad}
\end{equation}
\end{definition}

\begin{quote}
\textbf{Quoted Pavlov--Faddeev equivalence.}
A 2017 retrospective explicitly attributes to the 1972/1975
Pavlov--Faddeev paper the statement
\begin{equation}\label{eq:quoted-iff}
 \boxed{\qquad \mathrm{RH}\quad\Longleftrightarrow\quad
 \mathsf{PF}_{\mathrm{reg}}.\qquad}
\end{equation}
See \cite[\S3.2, formula (3.6), p.~1004]{TakhtajanEtAl2017}.
\end{quote}

This document records \eqref{eq:quoted-iff} as a \emph{source claim}, not as a
theorem rederived from the native definitions and proof in
\cite{PavlovFaddeev1975}.  A publication should verify the original
space, generator convention, hypotheses, and proof, and then establish their
equivalence with Hypotheses~\ref{hyp:energy}, \ref{hyp:LP},
\ref{hyp:minus-i}, \ref{hyp:inner}, and \ref{hyp:model}.

Under the unconditional lower bound proved conditionally on
Hypothesis~\ref{hyp:model} in Section~\ref{sec:easy},
$\mathsf{PF}_{\mathrm{reg}}$ is equivalent to
\begin{equation}\label{eq:PF-epsilon}
 \text{for every $\eps>0$ there is $C_{\eps}<\infty$ such that}
 \quad
 M(t)\leq C_{\eps}e^{(-1/4+\eps)t}
 \quad(t\geq0).
\end{equation}
The equivalence between the limsup upper bound and the family of
$\eps$-estimates in \eqref{eq:PF-epsilon} uses local boundedness:
$\sup_{0\leq t\leq T}\norm{Z(t)}<\infty$ for every $T<\infty$, and the
regularizer is fixed and bounded.  A finite initial interval is therefore
absorbed into $C_{\eps}$.  No estimate with a uniform constant at the boundary
exponent $-1/4$ is asserted.

\begin{problem}[Two separate handoff tasks]\label{prob:two-tasks}
\leavevmode
\begin{enumerate}[label=(T\arabic*)]
 \item \emph{Historical verification/reproof:} verify the original
       Pavlov--Faddeev statement and prove that its native objects are the
       objects defined here.  Reproving
       $\mathrm{RH}\Rightarrow\mathsf{PF}_{\mathrm{reg}}$ audits the quoted
       theorem but does not decide RH.
 \item \emph{Unconditional decision:} prove
       $\mathsf{PF}_{\mathrm{reg}}$ without assuming RH, or produce a rigorous
       failure of it that is tied to an off-critical-line zero.  By the quoted
       equivalence---conditional on verifying that its native objects satisfy
       Hypotheses~\ref{hyp:energy}, \ref{hyp:LP}, \ref{hyp:minus-i},
       \ref{hyp:inner}, and \ref{hyp:model}---this task is exactly the decision
       of RH.
\end{enumerate}
\end{problem}

\begin{warning}[Retain the limsup]
The family $Z(t)(B+iI)^{-1}$ is not itself a semigroup.  There is therefore no
submultiplicativity argument that automatically turns the limsup into a limit
or an infimum over $t>0$.
\end{warning}

\section{Conditional Hardy-space model and spectral input}\label{sec:model}

The elementary model below is exact once an intertwining theorem is supplied.
That physical-to-model theorem is isolated as Hypothesis~\ref{hyp:model};
nothing in this section silently promotes a scattering resonance to a
Hilbert-space eigenvalue.

Let
\[
 \C_{+}:=\{s\in\C:\Repart s>0\},
 \qquad
 H^{2}(\C_{+})
 :=\left\{f\text{ holomorphic}:
 \sup_{x>0}\frac1{2\pi}\int_{\R}|f(x+iy)|^{2}\,dy<\infty\right\}.
\]
Define the candidate reduced zeta model function
\[
 \Theta_{0}(s):=\frac{\xi(2s)}{\xi(-2s)}
               =\frac{\xi(2s)}{\xi(1+2s)},
 \qquad \Repart s>0,
\]
and isolate the function-theoretic input needed to call it inner.

\begin{hypothesis}[Reduced inner function]\label{hyp:inner}
The function $\Theta_{0}$ belongs to $H^{\infty}(\C_{+})$, has unimodular
boundary values almost everywhere on $i\R$, has zero mean type and no singular
inner factor, and its zeros in $\C_{+}$, with multiplicity, are exactly
$a=\rho/2$ over the nontrivial zeros of $\xi$.
\end{hypothesis}

Under Hypothesis~\ref{hyp:inner}, set
\[
 K_{0}:=H^{2}(\C_{+})\ominus\Theta_{0}H^{2}(\C_{+}),
 \qquad
 A_{0,t}:=P_{K_{0}}M_{e^{-ts}}\big|_{K_{0}}.
\]
Let $B_{0}$ generate $A_{0,t}^{*}$.  A zero $a$ of $\Theta_{0}$ gives a
reproducing-kernel eigenvector
\[
 A_{0,t}^{*}k_{a}=e^{-t\bar a}k_{a},
 \qquad
 B_{0}k_{a}=-\bar a\,k_{a}.
\]
For $a=\rho/2$, this is
\[
 \lambda_{\rho}=-\frac{\bar\rho}{2}.
\]
The standard modular Eisenstein scattering coefficient is
\[
 c(u)=
 \sqrt{\pi}\,
 \frac{\Gamma(u-\frac12)\zeta(2u-1)}
      {\Gamma(u)\zeta(2u)}.
\]
In the centered variable $u=1/2+s$, direct algebra gives
\begin{equation}\label{eq:c-vs-theta}
 c\!\left(\frac12+s\right)
 =\frac{s+\frac12}{s-\frac12}\,\Theta_{0}(s).
\end{equation}
Thus the physical meromorphic coefficient has an exceptional pole at
$s=1/2$, whereas $\Theta_{0}$ is the pole-cancelled inner factor.  The
half-plane Blaschke factor
\[
 b_{1/2}(s):=\frac{s-\frac12}{s+\frac12}
\]
cancels that pole, and \eqref{eq:c-vs-theta} says exactly
$b_{1/2}(s)c(1/2+s)=\Theta_{0}(s)$.  A causal characteristic function must be
derived after specifying whether it models $c$, $c^{-1}$, or a pole-cancelled
version.  In particular, the reciprocal of $b_{1/2}$ is meromorphic and is not
itself an inner factor on $\C_{+}$.

\begin{hypothesis}[Physical-to-model and resonance theorem]\label{hyp:model}
The following two assertions hold.
\begin{enumerate}[label=(M\arabic*)]
\item There are a finite-dimensional exceptional space
$\Kcal_{\mathrm{exc}}$, a semigroup $Z_{\mathrm{exc}}(t)$ on it, and a
boundedly invertible map
\[
 W:\Kcal\longrightarrow K_{0}\oplus\Kcal_{\mathrm{exc}}
\]
such that
\begin{equation}\label{eq:intertwine}
 WZ(t)W^{-1}=A_{0,t}^{*}\oplus Z_{\mathrm{exc}}(t).
\end{equation}
The exceptional generator has spectral bound at most $-1/2$.
\item The generator $B_{0}$ has meromorphic resolvent; every nontrivial zeta
zero $\rho$, with its algebraic multiplicity, gives the actual root space at
$-\bar\rho/2$; and
\begin{equation}\label{eq:spectrum}
 \sigma(B_{0})
 =
 \left\{-\frac{\bar\rho}{2}:\xi(\rho)=0\right\}.
\end{equation}
There is no additional boundary spectrum.
\end{enumerate}
\end{hypothesis}

The final sentence of Hypothesis~\ref{hyp:model} requires proof in any model
derivation.  For the reduced inner function it is supported by meromorphic
continuation across every finite point of $i\R$, absence of finite zero
accumulation, absence of a boundary singular inner factor, and zero mean type;
one must still cite the relevant model-generator spectral theorem.

An alternative historical parametrization is $(\rho-1)/2$.  The two multisets
agree after applying $\rho\mapsto1-\bar\rho$ to the zeta zeros.  Under
Hypothesis~\ref{hyp:model}, both conventions give
\begin{equation}\label{eq:spectral-bound}
 s(B):=\sup_{\lambda\in\sigma(B)}\Repart\lambda
 =-\frac12\inf_{\xi(\rho)=0}\Repart\rho.
\end{equation}

\begin{warning}[Effect of the intertwiner]
A bounded similarity preserves the exponential exponent because its condition
number contributes only $O(1/t)$ after taking logarithms.  Exact equality of
finite-time norms requires unitarity.  The finite exceptional summand cannot
change the target exponent because its spectral/growth exponent is at most
$-1/2$, strictly to the left of $-1/4$; this finite-dimensional claim must be
checked in the chosen full model.
\end{warning}

\section{Why the criterion implies RH}\label{sec:easy}

\begin{proposition}[Eigenvalue lower bound]\label{prop:eigen-lower}
Suppose $Bv=\lambda v$ for some nonzero $v\in\Dom B$ and
$\lambda\neq-i$.  Then
\[
 M(t)\geq\frac{e^{t\Repart\lambda}}{|\lambda+i|},
 \qquad
 \omega_{\mathrm{reg}}(B)\geq\Repart\lambda.
\]
\end{proposition}

\begin{proof}
Since $(B+iI)^{-1}v=(\lambda+i)^{-1}v$ and
$Z(t)v=e^{t\lambda}v$,
\[
 \frac{\norm{Z(t)(B+iI)^{-1}v}_{E}}{\norm{v}_{E}}
 =\frac{e^{t\Repart\lambda}}{|\lambda+i|}.
\]
Take the operator supremum and then the logarithmic limsup.
\end{proof}

Under Hypothesis~\ref{hyp:model}, applying
Proposition~\ref{prop:eigen-lower} to
$\lambda_{\rho}=-\bar\rho/2$ gives
\begin{equation}\label{eq:zero-lower}
 \omega_{\mathrm{reg}}(B)\geq-\frac12\Repart\rho
 \qquad\text{for every nontrivial zero $\rho$.}
\end{equation}
The zero set is nonempty and symmetric under
$\rho\mapsto1-\bar\rho$.  For each symmetric pair, at least one member has
real part at most $1/2$, so \eqref{eq:zero-lower} gives
$\omega_{\mathrm{reg}}(B)\geq-1/4$ unconditionally.  If
$\omega_{\mathrm{reg}}(B)\leq-1/4$, then \eqref{eq:zero-lower} excludes
every zero with $\Repart\rho<1/2$.  The symmetry
$\rho\mapsto1-\bar\rho$ then excludes every zero with
$\Repart\rho>1/2$.  Hence all nontrivial zeros lie on the critical line.

Thus the implication
\[
 \omega_{\mathrm{reg}}(B)\leq-\frac14
 \quad\Longrightarrow\quad\mathrm{RH}
\]
uses only the actual-eigenvector part of Hypothesis~\ref{hyp:model} and is not
the difficult analytic direction.  A pole of a meromorphically continued
scattering coefficient would not by itself suffice; the conversion to an
eigenvector in $\Dom B$ is essential.

\section{What an unconditional solution must establish}\label{sec:hard}

\begin{problem}[Primary solver handoff]\label{prob:primary}
First verify Hypotheses~\ref{hyp:energy}, \ref{hyp:LP},
\ref{hyp:minus-i}, \ref{hyp:inner}, and \ref{hyp:model} for the native
Pavlov--Faddeev semigroup.  Then, without assuming RH, prove
\begin{equation}\label{eq:hard}
 \forall\eps>0\ \exists C_{\eps}<\infty\ \forall t\geq0:
 \qquad
 \norm{Z(t)(B+iI)^{-1}}_{E\to E}
 \leq C_{\eps}e^{(-1/4+\eps)t}.
\end{equation}
Together with the lower bound in Section~\ref{sec:easy}, this proves
$\mathsf{PF}_{\mathrm{reg}}$, and Proposition~\ref{prop:eigen-lower} then
proves RH under Hypothesis~\ref{hyp:model}.
\end{problem}

Conversely, locating a zeta zero off the critical line and applying the
functional-equation symmetry produces a zero with real part below $1/2$;
that symmetric zero immediately violates \eqref{eq:hard} through
\eqref{eq:zero-lower} and disproves RH.  Merely showing
that the operator norm in \eqref{eq:hard} is too large for a non-spectral,
pseudospectral reason would instead refute the quoted implication
$\mathrm{RH}\Rightarrow\mathsf{PF}_{\mathrm{reg}}$; it would not, without the
verified Pavlov--Faddeev equivalence, disprove RH.  This distinction must be
maintained in any claimed negative solution.

If the goal is to audit the historical theorem rather than decide RH, the
conditional statement to reprove is
\begin{equation}\label{eq:conditional-hard}
 \mathrm{RH}\Longrightarrow\eqref{eq:hard}.
\end{equation}
That implication is nontrivial because the location of $\sigma(B)$ alone does
not control a nonnormal semigroup uniformly.

Plausible analytic frameworks for \eqref{eq:hard} include:
\begin{enumerate}[label=(R\arabic*)]
 \item a weighted $L^{2}$ resolvent estimate combined with Plancherel;
 \item a Hardy-space or Carleson-embedding estimate for the regularized model
       multiplier;
 \item a dyadic decomposition in spectral height with oscillatory
       cancellation; or
 \item a quantitative root-space synthesis estimate strong enough to control
       the graph-norm unit ball.
\end{enumerate}
Each route must exploit the explicit automorphic scattering/resolvent formula;
compactness, meromorphic continuation, or spectral location alone is
insufficient.

\begin{warning}[Heuristic obstruction to an absolute contour bound]
\label{warn:L1}
Do not adopt without proof the absolute operator-norm integrability of
\[
 \frac{d}{d\tau}
 \bigl(R(a+i\tau,B)R(\lambda_{0},B)\bigr)
 =
 -iR(a+i\tau,B)^{2}R(\lambda_{0},B)
\]
on a vertical line $a=-1/4+\eps$.  On a critical-line eigenvector with
eigenvalue $\lambda_n=-1/4+i\eta_n$, one obtains the lower bound
\[
 \norm{R(a+i\tau,B)^2R(\lambda_0,B)}
 \geq
 \frac{1}
 {(\eps^2+(\tau-\eta_n)^2)\abs{\lambda_0-\lambda_n}}.
\]
Together with sufficiently strong distinct-ordinate density and separation
estimates, these lower bounds would make the absolute
$L^{1}(\R;\Bcal(\Kcal))$ integral diverge over dyadic blocks.  This document
does not supply that arithmetic covering lemma, so divergence is recorded as
a serious obstruction rather than a proved theorem.  A more plausible contour
argument retains oscillation or uses an $L^{2}$/admissibility mechanism instead
of taking the absolute operator norm before integration.
\end{warning}

\section{Statements that must not be substituted for the target}
\label{sec:not-target}

\subsection{The unregularized operator norm}

Define
\[
 \norm{Z(t)}_{E\to E}
 :=\sup_{0\neq f\in\Kcal}
 \frac{\norm{Z(t)f}_{E}}{\norm{f}_{E}}.
\]
The assertion
\begin{equation}\label{eq:unregularized}
 \limsup_{t\to\infty}\frac1t\log\norm{Z(t)}_{E\to E}
 \leq-\frac14
\end{equation}
implies the regularized estimate, up to the fixed factor
$\norm{(B+iI)^{-1}}$, and is generally a strictly stronger demand.  It is not
the target in Definition~\ref{def:PF}.  A separate theorem for the reduced
scalar compressed-translation model gives $\norm{A_{0,t}}=1$ for all $t>0$,
so its unregularized exponent is $0$.  Exact finite-time norm one transfers to
the physical semigroup only through a unitary version of
Hypothesis~\ref{hyp:model}; a bounded similarity transfers only the exponent.
This model result does not conflict with the regularized criterion, because
the vectors nearly attaining the unregularized norm may depend on $t$ and
escape to high frequency; the resolvent suppresses precisely such
graph-norm-unbounded data.

The available 2002 edition of Lax's text displays the unregularized statement
as equation~(54), calls it sufficient for RH, and attributes necessity to
Pavlov--Faddeev~\cite{Lax2002}.  That later attribution is historically in
tension with the explicitly resolvent-regularized Pavlov--Faddeev criterion
quoted in \cite{TakhtajanEtAl2017}.  Until the original paper is checked, the
two should be reported as distinct source statements rather than harmonized by
assumption.

\subsection{A dense fixed-vector condition}

The Lax--Phillips 1980 exposition records the different sufficient condition
\begin{equation}\label{eq:dense-vector}
 \limsup_{t\to\infty}\frac1t\log\norm{Z(t)f}_{E}
 <-\frac14,
 \qquad f\in\Dcal,
\end{equation}
where $\Dcal$ is a dense no-standing-waves class
\cite[\S6, formula (6.19)]{LaxPhillips1980}.  Its quantifier order is
\[
 \exists\Dcal\subset\Kcal\text{ dense}\quad
 \forall f\in\Dcal\quad
 \limsup_{t\to\infty}\frac1t\log\norm{Z(t)f}_{E}<-\frac14.
\]
There is no supremum over $f$ after $t$ is chosen.  Consequently this condition
does not imply an operator-norm estimate without additional uniformity and
must not be substituted for \eqref{eq:PF-criterion}.  The strict inequality is
part of this separate sufficient condition; it is not the equality/upper-bound
form of the regularized operator criterion.

Under Hypothesis~\ref{hyp:model}, the strict dense-vector condition cannot in
fact hold on a dense set.  A known critical-line zero supplies an isolated
eigenvalue $\lambda$ with $\Repart\lambda=-1/4$ and nonzero finite-rank Riesz
projection $P_{\lambda}$.  If a vector has exponential rate strictly below
$-1/4$, commutation of $P_{\lambda}$ with $Z(t)$ forces
$P_{\lambda}f=0$.  Hence every such vector lies in the proper closed subspace
$\ker P_{\lambda}$, which cannot contain a dense class.  Thus the strict
historical display is a sufficient how-not-to-prove-RH condition, not an
RH-equivalent target in the model used here.

\section{Logical dependencies and prohibited shortcuts}

An attempted solution should identify which of the following inputs it uses.

\begin{enumerate}[label=(L\arabic*)]
 \item \textbf{Energy construction:} positivity of $E$ on the reduced
       no-standing-waves space and unitarity of $U(t)$ there.
 \item \textbf{Scattering construction:} the Lax--Phillips axioms for
       $\Dcal_{\pm}$ and the contraction-semigroup property of $Z(t)$.
 \item \textbf{Model theorem:} a precise unitary or boundedly invertible
       intertwiner between the physical semigroup and the reduced $K_{0}$
       realization plus its exceptional finite-dimensional summand, including
       the adjoint convention.
 \item \textbf{Resonance theorem:} meromorphic continuation of the generator
       resolvent and the zeta-zero/eigenvalue correspondence with algebraic
       multiplicities.
 \item \textbf{Uniform analytic estimate:} a resolvent, contour, Hardy-space,
       or admissibility estimate strong enough to prove \eqref{eq:hard}.
\end{enumerate}

The following shortcuts are invalid unless supplemented by new quantitative
arguments:
\begin{enumerate}[label=(S\arabic*)]
 \item $s(B)=-1/4$ does not imply the regularized or unregularized growth
       estimate for a nonnormal semigroup.
 \item Completeness of eigenvectors or generalized eigenvectors does not imply
       a uniform synthesis bound; a Riesz-basis or comparable quantitative
       statement would be needed.
 \item Compactness or meromorphicity of the resolvent does not control its norm
       on high vertical lines.
 \item A proof for every fixed $f$ does not permit exchanging
       $\sup_{f}$ with $\limsup_{t\to\infty}$.
 \item An approximate eigenvector used to prove $\norm{Z(t)}=1$ may depend on
       $t$; it is not a counterexample to the graph-regularized estimate.
 \item A Gearhart--Pr\"uss argument for the unregularized semigroup addresses a
       different norm and cannot simply discard the fixed right resolvent.
 \item Multiple zeta zeros produce root spaces and possible polynomial factors
       in $t$; multiplicities must be retained throughout contour and basis
       arguments.
 \item Replacing $Z$ by $Z^{*}$ preserves operator norms and real spectral
       bounds but not fixed-vector stability statements.  The convention must
       remain fixed through the proof.
\end{enumerate}

\section{Acceptance criteria for a claimed solution}

A proof or disproof handed back against this specification should contain all
of the following.

\begin{enumerate}[label=(A\arabic*)]
 \item It states whether it works on the physical energy space or on a model
       space and supplies the exact intertwining map if it transfers between
       them.
 \item It uses the positive reduced energy norm, not the raw indefinite form.
 \item It estimates exactly $Z(t)(B+iI)^{-1}$, or an explicitly proved
       equivalent fixed-resolvent/graph-norm quantity.
 \item It keeps the order of quantifiers
       $\forall\eps\,\exists C_{\eps}\,\forall t$.
 \item It proves operator-norm uniformity over the whole graph-norm unit ball.
 \item It treats the full scattering factor or proves that the omitted finite
       factor cannot change the logarithmic exponent.
 \item It tracks the generator/time/adjoint convention and the map from zeta
       zeros to resonances.
 \item It justifies every contour deformation, boundary limit, sum over poles,
       and interchange of limits or suprema.
 \item It handles zero multiplicities and generalized eigenvectors.
 \item It does not infer a growth estimate merely from spectral location.
\end{enumerate}

\section{Provenance and source status}

The regularized formula \eqref{eq:PF-criterion} is quoted as formula (3.6) in
the 2017 survey on Faddeev's scientific work, where it is explicitly attributed
to the Pavlov--Faddeev paper~\cite{TakhtajanEtAl2017,PavlovFaddeev1975}.  The
original 1972/1975 paper should be consulted for a publication-level check of
its exact native notation, hypotheses, and proof.  This document fixes a
mathematically unambiguous target conditional on the five labeled
identification hypotheses; it does not claim that freshly fixing notation
proves those identifications.

The Lax--Phillips monograph develops the automorphic energy and scattering
spaces~\cite{LaxPhillips1976}; the 1980 survey records the separate dense-vector
condition~\cite{LaxPhillips1980}.  Uetake gives an operator model for the
automorphic Lax--Phillips generator and scattering matrix~\cite{Uetake2007};
the exact theorem needed to establish Hypothesis~\ref{hyp:model}, with the
present sign and exceptional-factor conventions, must be cited explicitly in
any final proof.
General $C_{0}$-semigroup and resolvent facts used above can be found in
\cite{EngelNagel2000}.

\begin{thebibliography}{99}

\bibitem{PavlovFaddeev1975}
B.~S.~Pavlov and L.~D.~Faddeev,
\emph{Scattering theory and automorphic functions},
J. Soviet Math. \textbf{3} (1975), no.~4, 522--548;
English translation of Zap. Nauchn. Sem. LOMI \textbf{27} (1972), 161--193.

\bibitem{LaxPhillips1976}
P.~D.~Lax and R.~S.~Phillips,
\emph{Scattering Theory for Automorphic Functions},
Annals of Mathematics Studies, vol.~87,
Princeton University Press, Princeton, NJ, 1976.

\bibitem{LaxPhillips1980}
P.~D.~Lax and R.~S.~Phillips,
\emph{Scattering theory for automorphic functions},
Bull. Amer. Math. Soc. (N.S.) \textbf{2} (1980), no.~2, 261--295,
\href{https://doi.org/10.1090/S0273-0979-1980-14735-7}
{doi:10.1090/S0273-0979-1980-14735-7}.

\bibitem{Lax2002}
P.~D.~Lax,
\emph{Functional Analysis},
Wiley-Interscience, New York, 2002; see \S37.9, equation (54), p.~511.

\bibitem{TakhtajanEtAl2017}
L.~A.~Takhtajan, A.~Yu.~Alekseev, I.~Ya.~Aref'eva,
M.~A.~Semenov-Tian-Shansky, E.~K.~Sklyanin, F.~A.~Smirnov,
and S.~L.~Shatashvili,
\emph{Scientific heritage of L. D. Faddeev. Survey of papers},
Russian Math. Surveys \textbf{72} (2017), no.~6, 977--1081;
see \S3.2, formula (3.6), p.~1004,
\href{https://doi.org/10.1070/RM9799}{doi:10.1070/RM9799}.

\bibitem{Uetake2007}
Y.~Uetake,
\emph{The Lax--Phillips infinitesimal generator and the scattering matrix for
automorphic functions},
Ann. Polon. Math. \textbf{92} (2007), no.~2, 99--122,
\href{https://doi.org/10.4064/ap92-2-1}{doi:10.4064/ap92-2-1}.

\bibitem{EngelNagel2000}
K.-J.~Engel and R.~Nagel,
\emph{One-Parameter Semigroups for Linear Evolution Equations},
Graduate Texts in Mathematics, vol.~194,
Springer-Verlag, New York, 2000.

\end{thebibliography}

\end{document}
